% Paper 5, §7 — Safety evaluation, layer 1: adversarial robustness
% Anchors: kb/notes/alignment/safety-evaluation.md
%          kb/excerpts/{harmbench2024,chao2024-jailbreakbench,
%                       russinovich2024-crescendo,wei2023-jailbroken}.md

\section{Safety evaluation, layer 1: adversarial robustness}
\label{sec:adversarial-robustness}
\label{sec:safety-adversarial-robustness}

The agentic benchmarks of \cref{sec:agentic-benchmarks} ask what a
model \emph{can} do; the layer taken up here asks what a model
\emph{will} do under attack. The construct is no longer task
completion against a verifier $V(s_T, g)$; it is refusal under an
adversary whose explicit goal is to elicit a harmful behavior the
deployed system was trained to decline. The thesis of this section
is narrow: the field has built four canonical instruments for the
adversarial-robustness layer --- \citet{wei2023-jailbroken}'s
two-failure-mode taxonomy, \citet{harmbench2024}'s standardized
red-team suite, \citet{chao2024-jailbreakbench}'s reproducibility
benchmark, and \citet{russinovich2024-crescendo}'s multi-turn attack
--- and the four together do \emph{not} compose into a single
attack-success-rate (ASR) number that is comparable across papers.
A $5\%$ ASR on benchmark $X$ and a $30\%$ ASR on benchmark $Y$ for
the same model are not a contradiction; they are two measurements
of different constructs, instantiated against different behavior
sets, attack families, and judges. Before any \emph{within-paper}
ASR rank can be read as evidence of model robustness, the
\emph{across-paper} ASR-comparability problem has to be acknowledged
as the central methodological failure of the layer.

\subsection{Formal frame}
\label{sec:safety-formal}

A safety-evaluation instance%
\footnote{KB note: \texttt{kb/notes/alignment/safety-evaluation.md\#1-formal-framing}.}
is the tuple
\[
\mathcal{S} = (\mathcal{B}, \mathcal{A}, V, \pi),
\]
where $\mathcal{B}$ is the set of harmful behaviors to test (510 in
HarmBench, 100 in JailbreakBench), $\mathcal{A}$ is the attack family
(single-turn fixed prompt, optimization-based suffix, multi-turn
escalation, automated red-team), $V$ is the judge that decides
whether a generation realizes the behavior (rule-based, fine-tuned
classifier, LLM-as-judge, or human), and $\pi$ is the protocol
context (system prompt, decoding parameters, in-place defenses,
attack budget). The headline metric is the attack success rate
\citep[\S{}abstract]{chao2024-jailbreakbench}:
\[
\mathrm{ASR}(\theta, \mathcal{S}) =
\frac{|\{j \in J : V(\theta(\pi, a_j, b_j)) = \text{harmful}\}|}{|J|},
\]
the fraction of $(b, a)$ pairs $j \in J$ for which the model under
test, parameterized by $\theta$, produces an output the judge rates
harmful%
\footnote{KB note: \texttt{kb/notes/alignment/safety-evaluation.md\#1-formal-framing}, Eq.\,(1).}.
A perfectly aligned model has $\mathrm{ASR} = 0$ on every attack
family; a fully jailbroken model has $\mathrm{ASR} \approx 1$. The
load-bearing observation is that every term in the metric ---
$\mathcal{B}$, $\mathcal{A}$, $V$, $\pi$ --- varies across the four
canonical benchmarks of this section. ``ASR'' as a scalar is
under-specified without all four: the methodologically mature
reporting form is the tuple $(\mathcal{B}, \mathcal{A}, V, \pi,
\mathrm{ASR})$, not the number alone.

\begin{intuition}
ASR conflates three quantities under a single ratio: the model's
refusal robustness, the attack's strength, and the judge's
harshness. Two papers reporting ``GPT-4 ASR $= 30\%$'' and ``GPT-4
ASR $= 60\%$'' may both be correct under their respective
$(\mathcal{B}, \mathcal{A}, V, \pi)$ specifications. The $30\%
\to 60\%$ delta carries no information about the model unless the
specifications match. (Returning to canonical form: the ratio in
the ASR equation is well defined for any fixed $\mathcal{S}$;
across different $\mathcal{S}$ the ratios are not commensurable
quantities and arithmetic between them --- ranking, averaging,
delta-tracking --- is malformed.)
\end{intuition}

\subsection{The Wei taxonomy: why safety training fails}
\label{sec:safety-wei}

\citet{wei2023-jailbroken} supply the conceptual prior the
benchmarks instantiate. Their hypothesis is that successful
jailbreaks are not idiosyncratic exploits but realizations of two
structural failure modes of safety training%
\footnote{KB excerpt: \texttt{kb/excerpts/wei2023-jailbroken.md\#abstract}.}
\citep[\S{}abstract]{wei2023-jailbroken}:

\begin{quote}\itshape
We hypothesize two failure modes of safety training: \emph{competing
objectives} and \emph{mismatched generalization}. Competing
objectives arise when a model's capabilities and safety goals
conflict, while mismatched generalization occurs when safety
training fails to generalize to a domain for which capabilities
exist.
\end{quote}

Competing-objective attacks invoke a capability the model is
trained to deploy --- helpfulness, instruction-following, role-play
fidelity --- against the refusal training that would otherwise
trigger. Mismatched-generalization attacks exploit a gap between
the capability surface and the safety surface: capabilities cover
Base64, low-resource languages, ROT-13, structured role-play; safety
training was conducted in plain English direct-answer mode; the gap
is exploitable. The empirical finding is the load-bearing one
\citep[\S{}abstract]{wei2023-jailbroken}:

\begin{quote}\itshape
new attacks utilizing our failure modes succeed on every prompt in
a collection of unsafe requests from the models' red-teaming
evaluation sets and outperform existing ad hoc jailbreaks.
\end{quote}

Two things matter for the rest of the section. First, the
``every prompt'' result against GPT-4 and Claude v1.3 establishes
that extensively-red-teamed frontier models retain
\emph{conceptually predictable} vulnerabilities. The attacks are
not lucky; they instantiate the failure-mode framework. Second,
the policy upshot --- ``safety mechanisms should be as sophisticated
as the underlying model''
\citep[\S{}abstract]{wei2023-jailbroken} --- is the
\emph{safety-capability parity} prescription that subsequent
benchmarks evaluate against. HarmBench's $510 \times 18$ behavior-
attack matrix and JailbreakBench's leaderboard are operational
realizations of the test the Wei framework demands: cover the
capability surface broadly enough that mismatched-generalization
gaps surface; cover the attack surface broadly enough that
competing-objective constructions appear%
\footnote{KB note: \texttt{kb/notes/alignment/safety-evaluation.md\#2-1-adversarial-robustness-benchmarks}.}.

\subsection{HarmBench: the standardized red-team suite}
\label{sec:safety-harmbench}

\citet{harmbench2024} address the methodological mess in 2023
jailbreak literature by fixing the substrate against which attacks
and defenses are measured%
\footnote{KB excerpt: \texttt{kb/excerpts/harmbench2024.md\#abstract}.}
\citep[\S{}abstract]{harmbench2024}:

\begin{quote}\itshape
we introduce HarmBench, a standardized evaluation framework for
automated red teaming\ldots Using HarmBench, we conduct a
large-scale comparison of 18 red teaming methods and 33 target
LLMs and defenses, yielding novel insights. We also introduce a
highly efficient adversarial training method that greatly enhances
LLM robustness across a wide range of attacks, demonstrating how
HarmBench enables codevelopment of attacks and defenses.
\end{quote}

The framework's structural commitments
\citep[\S{}sec-structure]{harmbench2024}: \emph{510 carefully
curated behaviors} spanning four functional categories (standard
text, contextual, copyright, multimodal); \emph{18 red-team attack
methods} evaluated head-to-head (GCG, AutoDAN, PAIR, TAP, GBDA,
AutoPrompt, and successors\footnote{GCG and the
optimization-based-suffix lineage are deferred to AdvBench
\deepencite{zou2023-advbench, citation-pending} for full
treatment.}); \emph{33 target LLMs and defenses} measured against
the same suite; and a fine-tuned \emph{HarmBench-Judge} classifier
(Llama-2-13B family, fine-tuned on labeled harmful/non-harmful
generations) used as the standardized $V$.

The methodological move is structural rather than algorithmic:
HarmBench fixes $\mathcal{B}$, $\mathcal{A}$, and $V$ once, and
asks every attack and every defense to report numbers against the
same triple. The substrate is what enables co-development ---
``demonstrating how HarmBench enables codevelopment of attacks and
defenses''
\citep[\S{}abstract]{harmbench2024}. A new attack reports its
ASR on the HarmBench-510-behavior set under the
HarmBench-Judge against a fixed list of defenses; a new defense
reports its ASR-reduction on the same suite against a fixed list
of attacks. Cross-paper comparison becomes possible inside the
HarmBench frame in a way it was not in the 2023 literature, where
each paper picked its own behaviors, its own judge, and its own
target list.

The boundary the substrate does not cross is the one the next
benchmark addresses: HarmBench's adversarial-prompt corpus is a
research artifact, not a public reproducibility infrastructure
indexed against a leaderboard. The HarmBench-Judge classifier and
the behavior list are open; the per-attack adversarial prompts
generated by each of the 18 methods are not all republished as a
versioned artifact registry%
\footnote{KB note: \texttt{kb/notes/alignment/safety-evaluation.md\#2-1-adversarial-robustness-benchmarks}.}.

\subsection{JailbreakBench: the reproducibility leaderboard}
\label{sec:safety-jailbreakbench}

\citet{chao2024-jailbreakbench} are explicit that their
contribution is reproducibility infrastructure, not a new attack or
defense%
\footnote{KB excerpt: \texttt{kb/excerpts/chao2024-jailbreakbench.md\#abstract}.}
\citep[\S{}abstract]{chao2024-jailbreakbench}:

\begin{quote}\itshape
First, there is no clear standard of practice regarding
jailbreaking evaluation. Second, existing works compute costs and
success rates in incomparable ways. And third, numerous works are
not reproducible, as they withhold adversarial prompts, involve
closed-source code, or rely on evolving proprietary APIs. To
address these challenges, we introduce JailbreakBench, an
open-sourced benchmark with the following components: (1) an
evolving repository of state-of-the-art adversarial prompts, which
we refer to as jailbreak artifacts; (2) a jailbreaking dataset
comprising 100 behaviors\ldots\ which align with OpenAI's usage
policies; (3) a standardized evaluation framework that includes a
clearly defined threat model, system prompts, chat templates, and
scoring functions; and (4) a leaderboard that tracks the performance
of attacks and defenses for various LLMs.
\end{quote}

Four structural commitments. \emph{One hundred behaviors}, sourced
in part from prior work \citep{harmbench2024} and aligned with
OpenAI's usage policies. An \emph{evolving artifact repository}
that publishes the actual successful jailbreak prompts with
provenance, evaluation conditions, and judge verdict ---
operationally analogous to a CVE database for adversarial prompts.
A \emph{standardized $\pi$}: published threat model, system prompt,
chat template, and scoring function. And a \emph{public
leaderboard} that pins $(\mathcal{B}, V, \pi)$ and lets attacks and
defenses be ranked under a fixed denominator.

The complementarity with HarmBench is structural. HarmBench has
broader behavior coverage (510 vs.\ 100) and an attack-method
bake-off (18 attacks $\times$ 33 targets) inside one paper.
JailbreakBench has deeper artifact reproducibility and a public
leaderboard that admits new entries continuously%
\footnote{KB note: \texttt{kb/notes/alignment/safety-evaluation.md\#2-1-adversarial-robustness-benchmarks}.}.
HarmBench fixes the suite at publication time and runs a one-time
bake-off; JailbreakBench fixes the suite at publication time and
maintains it as a running leaderboard. Both fix the substrate
inside their own frame; neither makes its frame interoperable with
the other's. A score on HarmBench-510 under HarmBench-Judge is not
recoverable from a score on JailbreakBench-100 under
JailbreakBench's GPT-4-class judge, even for the same model.

\subsection{Crescendo: the multi-turn attack regime}
\label{sec:safety-crescendo}

\citet{russinovich2024-crescendo} document the failure mode the
single-turn benchmarks of \cref{sec:safety-harmbench,sec:safety-jailbreakbench}
cannot detect. The attack is a multi-turn escalation that begins
with a benign prompt and gradually steers the conversation toward
the harmful target%
\footnote{KB excerpt: \texttt{kb/excerpts/russinovich2024-crescendo.md\#abstract}.}
\citep[\S{}abstract]{russinovich2024-crescendo}:

\begin{quote}\itshape
Crescendo is a simple multi-turn jailbreak that interacts with the
model in a seemingly benign manner. It begins with a general
prompt or question about the task at hand and then gradually
escalates the dialogue by referencing the model's replies
progressively leading to a successful jailbreak.
\end{quote}

The mechanism is documented at the prompt-construction level. Open
with a general/benign question on the topic; ask a slightly more
specific follow-up; reference the model's own previous answer as
if it had committed; iterate until the model produces the harmful
behavior%
\footnote{KB note: \texttt{kb/notes/alignment/safety-evaluation.md\#2-1-adversarial-robustness-benchmarks}.}.
The exploitation hinges on the model over-weighting recent tokens
and text it has just produced itself: the model's prior answers
create a self-commitment gradient the attacker climbs, and because
each individual turn is benign, no single-turn refusal classifier
fires.

The headline number is the load-bearing one
\citep[\S{}abstract]{russinovich2024-crescendo}:

\begin{quote}\itshape
Crescendomation surpasses other state-of-the-art jailbreaking
techniques on the AdvBench subset dataset, achieving 29-61\%
higher performance on GPT-4 and 49-71\% on Gemini-Pro.
\end{quote}

The 29--71\% ASR \emph{uplift} over single-turn baselines ---
measured on the AdvBench subset \deepencite{zou2023-advbench,
citation-pending} under Crescendomation, the automated
implementation of the manual escalation strategy --- is not a new
ASR scalar; it is a delta against an existing ASR scalar. Under
the formal frame of \cref{sec:safety-formal}, Crescendomation
changes the attack family $\mathcal{A}$ from single-turn to
multi-turn while holding $(\mathcal{B}, V, \pi)$ approximately
fixed, and the resulting $\mathrm{ASR}$ is dramatically higher on
the same models against the same behaviors. The methodological
implication is structural: \emph{single-turn benchmarks
systematically underestimate the attack surface}. A model with
$\mathrm{ASR} < 5\%$ on JailbreakBench's single-turn behaviors may
register $\mathrm{ASR} > 30\%$ under Crescendo against an
overlapping behavior set.

\begin{analogy}
The Crescendo escalation curve is the social-engineering analog of
phishing under commitment escalation. Single-prompt jailbreaks are
the email asking ``please send password'' --- easy to refuse.
Crescendo is the conversation that opens with help on a benign
related task, builds rapport over several turns, and then asks the
prepared question --- with the model's own prior answers cited
back to it as commitments it has already made. The
\emph{commitment escalation} pattern documented in human
social-engineering literature predicts exactly this: each
acceded-to small request raises the conditional probability of
acceding to the next, larger one. Returning to canonical form: the
29--71\% ASR uplift is the quantitative measurement of how much
prior-self-reference shifts the refusal probability under the same
$(\mathcal{B}, V, \pi)$, with $\mathcal{A}$ as the only term that
varies.
\end{analogy}

The 2026 multi-turn picture: no equivalent of HarmBench or
JailbreakBench yet exists for multi-turn attacks. Standardizing
the regime requires three commitments the field has not yet made
--- a fixed turn budget, an attack-judge interaction protocol
that is not gameable by the attacker LLM, and a defense-reset
specification between turns%
\footnote{KB note: \texttt{kb/notes/alignment/safety-evaluation.md\#5-frontier}.}.
Until those land, multi-turn ASR is reported with paper-specific
$\mathcal{A}$ and $\pi$ in a way that compounds the comparability
problem the next subsection names directly.

\subsection{The ASR-comparability problem}
\label{sec:safety-asr-comparability}

The four benchmarks together expose the layer's central failure.
Each fixes its own $(\mathcal{B}, \mathcal{A}, V, \pi)$, and the
four fixed tuples are not interoperable. Three concrete failure
modes break ASR comparability across papers%
\footnote{KB note: \texttt{kb/notes/alignment/safety-evaluation.md\#1-formal-framing}.}.

\emph{Judge mis-specification.} JailbreakBench uses an LLM-as-judge
with published evaluation prompts; HarmBench uses the fine-tuned
HarmBench-Judge classifier; many 2023 papers use ad-hoc rule-based
``the model said anything other than I cannot'' verifiers; vendor
internal evals use proprietary judge stacks. Two judges scoring the
same generation as ``harmful'' or ``not harmful'' is not the same
event in two papers, even when the generation is identical
\citep[\S{}abstract]{chao2024-jailbreakbench}.

\emph{Behavior-set drift.} HarmBench's 510 behaviors span four
functional categories with explicit copyright and multimodal
sub-suites; JailbreakBench's 100 behaviors are aligned with OpenAI's
usage policies. The two distributions over harm categories are
not the same population, and a method strong on one set may
under-perform on the other for reasons that are properties of the
suite, not of the method%
\footnote{KB note: \texttt{kb/notes/alignment/safety-evaluation.md\#1-formal-framing}.}.

\emph{Defense / system-prompt drift.} Production deployments include
runtime monitoring, abuse detection, and account-level controls
that the benchmark $\pi$ does not model. Benchmark ASR systematically
\emph{overestimates} attack success against deployed systems
relative to standalone-API ASR, but the size of the gap is itself
deployment-specific and not measurable from outside the lab%
\footnote{KB note: \texttt{kb/notes/alignment/safety-evaluation.md\#1-formal-framing}.}.

The methodologically mature reporting form is the five-tuple
$(\mathcal{B}, \mathcal{A}, V, \pi, \mathrm{ASR})$, not a scalar.
JailbreakBench's leaderboard discipline --- fix the judge, the
behaviors, the threat model, vary only the attack and the model
under test --- is the cleanest current answer to the comparability
problem
\citep[\S{}abstract]{chao2024-jailbreakbench}, but is not yet
field-wide: model cards continue to publish ASR numbers without
all four substrate terms, and cross-benchmark comparisons in
informal venues continue to treat ASR as a model property.

\begin{contradiction}
\textbf{ASR is reported as a model attribute; it is not one.}
A single number ``GPT-4 ASR $= 30\%$'' carries no model-level
information without $(\mathcal{B}, \mathcal{A}, V, \pi)$.
\citet{harmbench2024} and \citet{chao2024-jailbreakbench} respond
with substrate standardization inside their own frames; neither
substrate is field-wide, and the two are not interoperable.
\citet{russinovich2024-crescendo} document a 29--71\% uplift when
the attack family $\mathcal{A}$ alone shifts from single-turn to
multi-turn, holding everything else approximately fixed --- which
demonstrates exactly how much of the headline ASR is determined by
substrate rather than model. As of 2026-Q2 there is no field-wide
$(\mathcal{B}, V, \pi)$ that the next benchmark can pin against.
The field publishes ASR scalars at a fictional level of precision
and treats them as if comparable; \cref{sec:contradictions} returns
to this as one of the layered-defense view's central open
contradictions.
\end{contradiction}

The closing observation is structural, not pessimistic. The
$\mathcal{S} = (\mathcal{B}, \mathcal{A}, V, \pi)$ frame is exactly
the right vocabulary for the comparability problem; the issue is
not measurement-theoretic but coordination-theoretic. HarmBench's
co-development substrate and JailbreakBench's leaderboard
discipline are partial standardizations whose boundary is the
boundary of each project's adoption. A field-wide pin --- a
HarmBench-or-equivalent standardized ASR with cross-benchmark
calibration set --- would close the contradiction; nothing about
the formal frame prevents it from being built. \Cref{sec:contradictions}
catalogs this as one of the layer's tractable open questions.

\subsection{Forward-link: the dangerous-capability layer}
\label{sec:safety-forward}

The construct in this section is \emph{can the model be
jailbroken into producing a harmful generation}, and the
$\mathcal{B}$ in $\mathcal{S}$ is a curated behavior list that
stands proxy for harm in deployment. \Cref{sec:dangerous-capability}
goes one layer deeper and asks the prior question: what dangerous
capabilities does the model possess in the first place. A
jailbroken refusal of an item the model could not realize would
be a paper artifact, not a deployment risk; a non-jailbroken
refusal of an item the model can realize via tool use, agentic
scaffolding, or out-of-distribution prompting is a deployment risk
the refusal eval cannot price. The bio-uplift, cyber-uplift,
persuasion, and autonomy/replication eval suites that anchor the
dangerous-capability layer measure the realizable construct
directly. The two layers compose: capability gates the harm
$\mathcal{B}$ contains, refusal robustness gates the model's
willingness to act on it. \Cref{sec:sycophancy,sec:scheming,sec:alignment-faking}
shift one layer further, from \emph{what an attacker can elicit}
to \emph{what the model itself does under conflicting incentives}
--- the alignment-threat-detection regime, where the construct is
no longer adversarial robustness but the model's own behavior under
ordinary deployment, with no attacker present.
